% Options for packages loaded elsewhere
% Options for packages loaded elsewhere
\PassOptionsToPackage{unicode}{hyperref}
\PassOptionsToPackage{hyphens}{url}
%
\documentclass[
  11pt,
]{article}
\usepackage{xcolor}
\usepackage[margin=1in]{geometry}
\usepackage{amsmath,amssymb}
\setcounter{secnumdepth}{-\maxdimen} % remove section numbering
\usepackage{iftex}
\ifPDFTeX
  \usepackage[T1]{fontenc}
  \usepackage[utf8]{inputenc}
  \usepackage{textcomp} % provide euro and other symbols
\else % if luatex or xetex
  \usepackage{unicode-math} % this also loads fontspec
  \defaultfontfeatures{Scale=MatchLowercase}
  \defaultfontfeatures[\rmfamily]{Ligatures=TeX,Scale=1}
\fi
\usepackage{lmodern}
\ifPDFTeX\else
  % xetex/luatex font selection
\fi
% Use upquote if available, for straight quotes in verbatim environments
\IfFileExists{upquote.sty}{\usepackage{upquote}}{}
\IfFileExists{microtype.sty}{% use microtype if available
  \usepackage[]{microtype}
  \UseMicrotypeSet[protrusion]{basicmath} % disable protrusion for tt fonts
}{}
\makeatletter
\@ifundefined{KOMAClassName}{% if non-KOMA class
  \IfFileExists{parskip.sty}{%
    \usepackage{parskip}
  }{% else
    \setlength{\parindent}{0pt}
    \setlength{\parskip}{6pt plus 2pt minus 1pt}}
}{% if KOMA class
  \KOMAoptions{parskip=half}}
\makeatother
% Make \paragraph and \subparagraph free-standing
\makeatletter
\ifx\paragraph\undefined\else
  \let\oldparagraph\paragraph
  \renewcommand{\paragraph}{
    \@ifstar
      \xxxParagraphStar
      \xxxParagraphNoStar
  }
  \newcommand{\xxxParagraphStar}[1]{\oldparagraph*{#1}\mbox{}}
  \newcommand{\xxxParagraphNoStar}[1]{\oldparagraph{#1}\mbox{}}
\fi
\ifx\subparagraph\undefined\else
  \let\oldsubparagraph\subparagraph
  \renewcommand{\subparagraph}{
    \@ifstar
      \xxxSubParagraphStar
      \xxxSubParagraphNoStar
  }
  \newcommand{\xxxSubParagraphStar}[1]{\oldsubparagraph*{#1}\mbox{}}
  \newcommand{\xxxSubParagraphNoStar}[1]{\oldsubparagraph{#1}\mbox{}}
\fi
\makeatother


\usepackage{longtable,booktabs,array}
\usepackage{calc} % for calculating minipage widths
% Correct order of tables after \paragraph or \subparagraph
\usepackage{etoolbox}
\makeatletter
\patchcmd\longtable{\par}{\if@noskipsec\mbox{}\fi\par}{}{}
\makeatother
% Allow footnotes in longtable head/foot
\IfFileExists{footnotehyper.sty}{\usepackage{footnotehyper}}{\usepackage{footnote}}
\makesavenoteenv{longtable}
\usepackage{graphicx}
\makeatletter
\newsavebox\pandoc@box
\newcommand*\pandocbounded[1]{% scales image to fit in text height/width
  \sbox\pandoc@box{#1}%
  \Gscale@div\@tempa{\textheight}{\dimexpr\ht\pandoc@box+\dp\pandoc@box\relax}%
  \Gscale@div\@tempb{\linewidth}{\wd\pandoc@box}%
  \ifdim\@tempb\p@<\@tempa\p@\let\@tempa\@tempb\fi% select the smaller of both
  \ifdim\@tempa\p@<\p@\scalebox{\@tempa}{\usebox\pandoc@box}%
  \else\usebox{\pandoc@box}%
  \fi%
}
% Set default figure placement to htbp
\def\fps@figure{htbp}
\makeatother





\setlength{\emergencystretch}{3em} % prevent overfull lines

\providecommand{\tightlist}{%
  \setlength{\itemsep}{0pt}\setlength{\parskip}{0pt}}



 


\makeatletter
\@ifpackageloaded{caption}{}{\usepackage{caption}}
\AtBeginDocument{%
\ifdefined\contentsname
  \renewcommand*\contentsname{Table of contents}
\else
  \newcommand\contentsname{Table of contents}
\fi
\ifdefined\listfigurename
  \renewcommand*\listfigurename{List of Figures}
\else
  \newcommand\listfigurename{List of Figures}
\fi
\ifdefined\listtablename
  \renewcommand*\listtablename{List of Tables}
\else
  \newcommand\listtablename{List of Tables}
\fi
\ifdefined\figurename
  \renewcommand*\figurename{Figure}
\else
  \newcommand\figurename{Figure}
\fi
\ifdefined\tablename
  \renewcommand*\tablename{Table}
\else
  \newcommand\tablename{Table}
\fi
}
\@ifpackageloaded{float}{}{\usepackage{float}}
\floatstyle{ruled}
\@ifundefined{c@chapter}{\newfloat{codelisting}{h}{lop}}{\newfloat{codelisting}{h}{lop}[chapter]}
\floatname{codelisting}{Listing}
\newcommand*\listoflistings{\listof{codelisting}{List of Listings}}
\makeatother
\makeatletter
\makeatother
\makeatletter
\@ifpackageloaded{caption}{}{\usepackage{caption}}
\@ifpackageloaded{subcaption}{}{\usepackage{subcaption}}
\makeatother
\usepackage{bookmark}
\IfFileExists{xurl.sty}{\usepackage{xurl}}{} % add URL line breaks if available
\urlstyle{same}
\hypersetup{
  pdftitle={Referee Report (Second Round)},
  pdfauthor={Anonymous Referee},
  hidelinks,
  pdfcreator={LaTeX via pandoc}}


\title{Referee Report (Second Round)}
\author{Anonymous Referee}
\date{2026-07-24}
\begin{document}
\maketitle


\textbf{Date:} 2026-07-24 \textbf{Manuscript ID:} JEEM-D-26-00040R1
\textbf{Reviewer:} Anonymous Referee (Reviewer 2)

\subsection{Recommendation}\label{recommendation}

\textbf{Category: Minor Revision}

\subsection{Confidential Comments to
Editor}\label{confidential-comments-to-editor}

\subsubsection{Assessment of the
Revision}\label{assessment-of-the-revision}

This is a serious and largely successful revision, and I cross-checked
the response letter's claims against the revised manuscript; the
responses are accurate. The four central issues from the first round are
largely resolved. The new event study and the reframing of the
contribution as a reduced-form effect of the clean-air governance regime
settle the timing and interpretation concern; the rewritten multitask
model now maps one-to-one to the empirical specifications (which also
resolves Referee 1's model concerns); and the promotion coding, sample
restrictions, and IV diagnostics are documented and stress-tested. The
one partially outstanding area is validation of the text-based effort
measure.

\subsubsection{Remaining Required and Suggested
Concerns}\label{remaining-required-and-suggested-concerns}

My remaining points are limited. Two could in principle move results
rather than wording: the unreported lower-order interaction terms in the
information-gains specification (Table 3), on which the paper's
remaining mechanism claim rests, and a direct human-audit validation of
the AI-based text classifier, which the authors substituted with coarser
proxies. The rest are presentational: calibrating the abstract's
mechanism claim to the paper's own hedged language, reconciling an
OLS-IV sign contrast in the provincial-leader appendix, and documenting
an estimator label. The paper is close to being publishable at JEEM, and
I expect a quick final round.

\subsection{Comments to Authors}\label{comments-to-authors}

\subsubsection{Summary}\label{summary}

Thank you for the careful and substantive revision. The manuscript has
improved considerably. The conceptual framework is now a coherent
multitask principal-agent model whose first-order condition maps
directly to the empirical specifications, and the contribution has been
appropriately reframed as a reduced-form effect of the clean-air
governance regime. The promotion coding is documented and stress-tested,
and the new event-study, information-gains, local-projection, and
provincial-leader analyses respond directly to the first-round comments.
I verified the response letter's claims against the revised manuscript
and found the responses accurate.

My remaining points are limited. One substantive request from the first
round is still outstanding (validation of the AI-based text classifier),
and the new analyses introduce a small number of reporting and
interpretation issues. None of these requires new identification
strategies or major data work. I list them below as required points,
followed by optional suggestions.

\subsubsection{Assessment of Responses to First-Round
Comments}\label{assessment-of-responses-to-first-round-comments}

\textbf{3.1 (Core interpretation).} Largely addressed. The event study
(Figure 3) is presented candidly: pre-adoption coefficients are
imprecise without a systematic trend, and the post-adoption gradient
strengthens gradually, becoming significant after roughly three event
years. The reframing as a reduced-form effect of the
monitoring-and-disclosure regime is carried through the title, the
results text (p.~23), and the conclusion, which now states explicitly
that the paper does not directly test the weight-change channels. The
new information-gains test (Table 3) is a constructive addition and does
provide evidence consistent with an observability channel. Several
issues with this test and its framing remain. See Required Point R1.
\textbf{3.2 (Model-empirics mapping).} Fully resolved. The revised model
(Section 3) includes the multitask promotion schedule, CARA utility over
the realized career payoff with the risk premium carrying signal
variance, a well-motivated objective for the higher-level government,
and a first-order condition that is exactly the net-promotion-return
condition requested. Proposition 1's four-primitive comparative statics
give the empirical results an honest interpretation. \textbf{3.3
(Promotion coding and sample restrictions).} Resolved. The coding rule
for lateral transfers to sub-provincial cities and provincial capitals
is now explicit (Section 5.1 and footnote 6), demoted/disciplined
officials are included in the baseline (footnote 5; Table 1), and Table
7 columns (3) and (5) show the results are robust to excluding
sub-provincial cities and demoted/disciplined officials, respectively.
\textbf{3.4 (Instrument exclusion and diagnostics).} Substantially
addressed. The expanded discussion of the exclusion restriction and the
weather-control robustness check (Table 7, column 4) are welcome. Two
residual points on characterization and interpretation remain. See
Required Point R3. \textbf{3.5 (Effort measurement validation).}
Partially addressed. The category-level adoption trends (Appendix Figure
2), the released classification code, and the waste-gas investment proxy
(Appendix Tables 4 and 6) are useful steps. However, the core
first-round request - a direct validation of the AI-based classification
against human coding - has not been carried out, and the category-level
trends themselves show that several categories break exactly at
2013-2014, which keeps the reporting-template concern alive. See
Required Point R2. \textbf{3.6 (Fog vs.~smog interpretation).} Resolved.
The interpretation is now presented as one possibility.
\textbf{Secondary comments (4.1-4.6).} All substantively addressed. The
pre-period narrative is now carefully worded (the largest pre-period
t-statistic across Table 7 corresponds to p approx. 0.31), the log(1+Y)
transformation has been replaced and PPML added, and local projections
appear in Appendix Figure 4. The provincial-leader extension (Section 9,
Appendix Table 7) includes a frank power discussion, the heterogeneity
analysis (Table 4) is informative, and the documentation is much
expanded. Three small residuals from these items appear under Required
Points R4-R5 and one minor suggestion.

\subsection{Major Concerns}\label{major-concerns}

\textbf{R1. Complete Table 3's reporting and calibrate the mechanism
claim to it.} The paper's remaining mechanism claim rests on Table 3.
Four related fixes, only the first involving any estimation, would let
it bear that weight. (i) Report the constituent two-way terms log(PM2.5)
x API correlation and Post x API correlation. If these terms are omitted
from the specification, the triple interaction conflates the
information-gains effect with time-invariant differences in the
pollution-promotion gradient across high- and low-correlation cities and
with differential post-period level shifts; if they are included but
unreported, please report them or state their inclusion in the table
notes. (ii) Clarify the direction of the interaction variable. The text
narrates that ``larger information gains magnify the difference,'' but
the regressor is the API \emph{correlation}, which moves inversely with
the information gain, so the expected and estimated triple coefficient
is \emph{positive}. Either define an explicit information-gain variable
(for example, one minus the correlation) or add a clarifying sentence so
readers do not misread the sign. Note also that the implied post-shift
gradient change remains substantially negative even at the highest
observed correlation levels, which suggests that non-information
channels carry much of the effect, consistent with the paper's own
reduced-form framing. (iii) Add one sentence acknowledging that low
pre-period API correlation may proxy for pre-period misreporting
intensity, which could correlate with post-2013 enforcement exposure
through channels beyond information per se (larger revealed pollution
gaps, tighter targets, more central attention); the test is consistent
with the observability mechanism but does not isolate it from this
exposure channel. (iv) Align the abstract with Section 7.1's own
language. ``We show that at least part of this change occurred through
an observability-and-incentives mechanism'' overstates what the design
isolates; ``We provide evidence consistent with an
observability-and-incentives mechanism'' would match it. Similarly,
``mayors who reduced PM2.5 concentrations were no more likely to be
promoted before the regime shift'' reads as a precise zero, while the IV
pre-period estimate is an imprecise -0.275 (SE 0.298); ``not detectably
more likely'' (or similar) would be accurate. \textbf{R2. Validate the
text classifier directly.} The first-round request for a human-audit
validation of the AI-based classification remains open, and it is the
one remaining item that involves real (though modest) work: draw a
random audit sample of work-report passages, have them coded by human
coders blind to the AI output, and report agreement statistics (for
example, precision and recall by category). This matters because
Appendix Figure 2 shows that several categories - clean heating and coal
substitution, heavy-pollution emergency response, dust control - rise or
jump precisely at 2013-2014, and emergency-response plans were mandated
by the 2013 Action Plan. For these categories, reporting-template
changes and genuine effort changes are observationally similar in
keyword counts. A robustness check excluding Action-Plan-mandated
categories (or at a minimum tempering ``the adoption frequency of most
policy categories remained relatively stable'' to acknowledge the
exceptions) would substantially strengthen Section 7.3. \textbf{R3.
Discuss the weather-control sensitivity and characterize Appendix Table
2 accurately.} Two related points. First, Appendix Table 2 shows that
inversion strength predicts visibility (2.852, SE 1.555) and dew point
(1.349, SE 0.786), both significant at the 10 percent level; only
minimum temperature is a clean null. The current characterization
(``limited evidence \ldots{} significantly predicted'') understates two
marginally significant coefficients out of three. Second, adding the
weather controls moves the key interaction from -0.313 to -0.569 (Table
7, column 4): the sign and significance are stable, but the magnitude
nearly doubles. A brief discussion of why the estimate is sensitive in
magnitude to these controls (and which estimate the authors regard as
preferred, and why) would make the robustness claim complete. As
written, ``remain unchanged'' describes sign and significance only.
\textbf{R4. Document the estimator in Appendix Table 5.} The response
letter states that Appendix Table 5 uses PPML with fixed effects, and
Appendix Table 3's notes state Poisson quasi-maximum likelihood
explicitly, but Appendix Table 5's notes do not name the estimator.
Please state it in the notes (and report the same fit diagnostics as the
neighboring tables). \textbf{R5. Reconcile the OLS and IV results for
provincial leaders.} In Appendix Table 7, the OLS specification with
province controls (column 2) shows a positive and statistically
significant interaction (0.282, SE 0.129), while the IV interaction
(column 3) is negative and insignificant. The text discusses only the IV
null. One sentence acknowledging the OLS-IV contrast and explaining why
the IV estimate is preferred (and what the contrast implies, given only
22 clusters) would preempt confusion.

\subsection{Minor Concerns}\label{minor-concerns}

Section 3: the career-reward slope on growth (delta\_t) is taken as
given while the pollution slope (gamma\_t) is optimized. A one-sentence
justification for this asymmetry (for example, the growth slope
reflecting a long-standing institutional legacy) would tighten the
exposition. Appendix Figure 4: for the promotion outcome, the horizon-3
confidence interval includes zero and horizon 1 is borderline; ``robust
to alternative outcome horizons'' is slightly generous for that outcome
(the effort outcomes are robust throughout). Alternative promotion
definitions (for example, strict upward-only, or a multinomial treatment
separating adverse exits) remain a useful, optional robustness exercise
now that demoted/disciplined officials are in the baseline.

\subsubsection{Closing}\label{closing}

The revision substantially improves the paper and engages seriously with
every first-round comment. The required points above are bounded: one
modest validation exercise, one specification-reporting fix, and several
presentational calibrations. I expect they can be addressed quickly.




\end{document}
